\par\Needspace{6\baselineskip}\medskip
\hypertarget{native-120-article089}{}\pdfbookmark[2]{89.}{article-089}
\begin{center}89.\end{center}
Duplex iam negotium nobis incumbit, primum, ut ex aequatione transcendente 12, quoniam solutionem directam non admittit, incognitam $g$ quam commodissime eruamus; secundum, ut ex angulo $g$ invento elementa ipsa deducamus. Quae antequam adeamus, transformationem quandam attingemus, \SourceLeaf{121}cuius adiumento calculus quantitatis auxiliaris $l$ vel $L$ expeditius absolvitur, insuperque plures formulae post evolvendae ad formam elegantiorem reducuntur.

Introducendo scilicet angulum auxiliarem $\omega$ per formulam
\[
\BareRoot[4]{\dfrac{r'}{r}}=\operatorname{tang}(45^0+\omega)
\]
determinandum, fit
\[
\begin{aligned}
\BareRoot{\dfrac{r'}{r}}+\BareRoot{\dfrac{r}{r'}}
&=2+\{\operatorname{tang}(45^0+\omega)-\operatorname{cotang}(45^0+\omega)\}^2\\[4pt]
&=2+4\operatorname{tang}2\omega^2;
\end{aligned}
\]
unde habetur
\[
\begin{aligned}
l&=\frac{\sin\tfrac12f^2}{\cos f}+\frac{\operatorname{tang}2\omega^2}{\cos f}\mathpunct{\raisebox{.45ex}{,}}\\[6pt]
L&=-\frac{\sin\tfrac12f^2}{\cos f}-\frac{\operatorname{tang}2\omega^2}{\cos f}\mathpunct{\cdot}
\end{aligned}
\]
